\begin{officialform}[BulletinEN]\begin{english}
\FormHeading
\FormLine{Project code: \ProjectCode}
\FormLine{Project title: \ProjectTitleEnglish}
\FormLine{Duration: \ProjectPeriodEnglish}
\FormLine{Budget: \ProjectBudgetEnglish}
\FormLine{Coordinator: \PrincipalName\ (\PrincipalPhone)}
\FormLine{Members: \ProjectMembers}
\FormLine{Supervisor: \SupervisorName\ (\SupervisorEmail)}
\FormLine{Project rationale: Universities and organizations increasingly need task management and collaboration tools. A shared service must identify each tenant correctly, protect its data, and maintain service quality while tenants use common infrastructure.}
\FormLine{Objectives: To study, design, and develop a multi-tenant Kanban task management application; deploy it as an Internet-accessible service; and evaluate its data isolation, performance, and runtime customization.}
\FormLine{Methodology: Literature review and analysis; requirements-driven architecture design; prototype implementation; and controlled experiments covering identity and context integrity, data isolation, performance, and tenant-specific customization.}
\FormLine{Project activities: Review multi-tenant architectures; analyze requirements; design tenant identification, data isolation, rate limiting, caching, and customization mechanisms; build and deploy the Kanban application; conduct experiments and evaluate results.}
\FormLine{Research results: The project developed TenantFlow using a Level 3 Pool architecture and deployed it on AWS EC2. The Chapter 5 summary reports no cross-tenant data leaks in 600 fault-injection cases and no context contamination across 1,500 alternating tenant requests. With 80 workload virtual users (up to 82 including test users), p95 latency for regular tenants was approximately 242--277 ms. All 19 tested extended-field card creations succeeded. Tail latency at p99 remained approximately 1.99--2.45 seconds and requires further optimization.}
\FormLine{Research contribution: The study connects a multi-tenant architecture framework with a concrete Kanban application and evaluates identity controls, database row-level security, rate limiting, two-level caching, and runtime customization together in a deployed prototype.}
\FormLine{Potential applications: The prototype and findings can inform task management and collaboration services for universities, small and medium-sized businesses, and startup teams.}
\end{english}\end{officialform}
